% Shared by both versions: the dedicated credit leaf. Source: spec
% 7.5.3 and 10.3 (0.11.0). The heading and label are the includer's:
% a run-in paragraph of the voucher section (short) or of the receive
% section (full-receive.tex), so the text starts without a blank line.
A dedicated credit leaf is a source leaf that the issuer owns and whose
exit key is one recipient's, in a tree whose taptree carries
\tr{CLAIM}$^{x}$, \tr{UPDATE}$^{x}$ and the exact-payout exit, carries
\tr{REWRITE}$_h^{x}$ in place of \tr{REWRITE}$_h$ (or no rewrite at all),
and omits \tr{SPLIT}. \tr{CLAIM}$^{x}$ is \tr{CLAIM} with one further
signature over the transaction, by the source leaf's exit key: 7 bytes
of script and 64 of witness in FULL, the only tier that builds it. The
recipient's own redemptions carry that signature at no cost, and the
issuer cannot carve the leaf at all, not even with a voucher it signs to
itself. Over-issuance then has no second victim, since every voucher
pays the same recipient, who can watch the leaf's amount. Voiding needs
the leaf to change, and every transition that could change it needs the
exit key: \tr{CLAIM}$^{x}$, \tr{UPDATE}$^{x}$ and \tr{REWRITE}$_h^{x}$
by rule, while the exact exit and an eviction pay the whole balance to
that key and so void nothing the recipient is owed. The balance can only grow or pay its recipient, whose remaining
exposure is the issuer's liveness: a provider that stops issuing.

Three details keep it that way. The recipient checks from the taptree
root $m$ that no plain \tr{REWRITE}$_h$ is present, since it needs only
the owner's and the operator's signatures, both the issuer's. The issuer
funds the leaf only with credit it has issued or is about to, since any
excess is the recipient's on every path the issuer can take alone, and
tops it up by an amount-only increase under \tr{REWRITE}$_h^{x}$, which
needs only the operator's signature. And the leaf has a finite lease
with the operator as evictor, not a $\NUMS$ evictor, which every carve
would inherit, blocking the sweep of whatever cohort the claim leaf lands
in. The credit is one-way: if the recipient could also spend from it, an
older, larger voucher would be a stale state needing a window to
supersede, as in a bidirectional channel. To spend, the recipient
redeems into a channel leaf (\S\ref{sec:channels}).
